\section*{Resumen ejecutivo}
\addcontentsline{toc}{section}{Resumen ejecutivo}

El presente trabajo analiza tres series temporales mediante modelos de la familia ARIMA y
un modelo de vectores autorregresivos. Dos de las series provienen de la infraestructura
productiva de una plataforma comercial y registran, con frecuencia horaria, el volumen de
solicitudes atendidas por un balanceador de carga y por una interfaz de punto de venta. La
tercera proviene de la telemetría de vuelo de una aeronave no tripulada de ala fija y
registra su altitud con frecuencia de un segundo.

Las dos series de tráfico exhiben estacionalidad diaria con período de veinticuatro horas y
quedan representadas por modelos SARIMA. La serie de altitud resulta un proceso integrado sin
componente estacional y queda representada por un modelo ARIMA$(1,1,3)$, cuyos parámetros son
significativos tanto individual como conjuntamente, y cuyos residuos no se distinguen del
ruido blanco.

El modelo de vectores autorregresivos se construye sobre las dos series de tráfico, únicas que
comparten calendario. La especificación adoptada emplea dos rezagos, respaldada por los
criterios bayesiano y de Hannan-Quinn. La prueba de causalidad de Granger arroja un resultado
unidireccional: los rezagos de la serie de punto de venta contribuyen a explicar la serie del
balanceador, mientras que la relación inversa no alcanza significatividad.

Los principales límites del trabajo son el tamaño muestral de cada serie, el solapamiento de
apenas trece días entre las dos series de tráfico, y la imposibilidad de incorporar la serie
de altitud al modelo multivariado por corresponder a un calendario disjunto.
